\chapter*{คำนำ}
\addcontentsline{toc}{chapter}{คำนำ}

ตำราวิชาการเรื่อง \textbf{หลักสูตรเกษตรดิจิทัลและปัญญาประดิษฐ์ฝังตัว สวนผลไม้มูลค่าสูงและการเพาะเลี้ยงสัตว์น้ำชายฝั่ง ภาคตะวันออก} เล่มนี้ เรียบเรียงขึ้นเพื่อใช้เป็นตำราและเอกสารประกอบการฝึกอบรมขั้นสูงในโครงการพัฒนาบัณฑิตพันธุ์ใหม่และกำลังคนสมรรถนะสูงเพื่อตอบโจทย์ภาคการผลิตตามนโยบายการปฏิรูปการอุดมศึกษาไทย โดยมุ่งเน้นการบูรณาการองค์ความรู้ทางฟิสิกส์ประยุกต์ เคมีเชิงฟิสิกส์ วิทยาการคอมพิวเตอร์ และเทคโนโลยีปัญญาประดิษฐ์ฝังตัว (Embedded AI และ Edge Computing) เข้ากับบริบทการแก้ปัญหาจริงในแปลงเกษตรกรรมและฟาร์มเพาะเลี้ยงสัตว์น้ำเศรษฐกิจของภาคตะวันออก

พื้นที่ภาคตะวันออกของประเทศไทย โดยเฉพาะจังหวัดจันทบุรี ระยอง และตราด ถือเป็นศูนย์กลางทางเศรษฐกิจการเกษตรที่สำคัญอย่างยิ่ง ทั้งในด้านการผลิตผลไม้เมืองร้อนมูลค่าสูง เช่น ทุเรียน มังคุด และเงาะ รวมทั้งเป็นแหล่งเพาะเลี้ยงสัตว์น้ำชายฝั่งที่สำคัญ เช่น กุ้งขาวแวนนาไมและปลากะพงขาว อย่างไรก็ดี เกษตรกรและผู้ประกอบการต้องเผชิญกับความท้าทายเชิงโครงสร้าง ทั้งปัญหาการเปลี่ยนแปลงสภาพภูมิอากาศ ความผันผวนของต้นทุนพลังงานไฟฟ้า แรงงานขาดแคลน และข้อกำหนดด้านมาตรฐานคุณภาพผลผลิตส่งออกที่เข้มงวด การทำการเกษตรแบบพึ่งพาประสบการณ์ดั้งเดิมเพียงอย่างเดียวจึงไม่เพียงพอต่อการรักษาขีดความสามารถในการแข่งขัน

คณะผู้เขียนจึงได้สังเคราะห์กรอบการทำงานเชิงระบบโดยแบ่งออกเป็นสองกิจกรรมหลัก ได้แก่
\begin{enumerate}
    \item \textbf{กิจกรรมที่ 1 นวัตกร AIoT สวนผลไม้มูลค่าสูง} มุ่งเน้นการตรวจวัดความสมบูรณ์ของดินและสภาพอากาศด้วยโพรบ 7-in-1 การประเมินค่าแรงดึงระเหยน้ำของอากาศ (VPD) การคำนวณการใช้น้ำของพืช ($ET_c$) ตามระยะฟีโนโลยี การควบคุมระบบไฮดรอลิกส์ด้วยวาล์วไฟฟ้า 24VAC ป้องกันปรากฏการณ์ Water Hammer การประเมินความสุกแก่ของทุเรียนด้วยคอมพิวเตอร์วิทัศน์สีหนามร่องพูและการสั่นพ้องของคลื่นเสียง และระบบคัดเกรดผลไม้อัตโนมัติบนสายพานด้วยโมเดล YOLO
    \item \textbf{กิจกรรมที่ 2 นวัตกร AIoT ฟาร์มเพาะเลี้ยงสัตว์น้ำชายฝั่ง} มุ่งเน้นการตรวจวัดคุณภาพน้ำแบบเรียลไทม์ด้วยโพรบวัดออกซิเจนละลายเชิงแสง (Optical DO) และความเค็ม การคำนวณอุณหพลศาสตร์การละลายของออกซิเจน (Benson-Krause Model) การควบคุมมอเตอร์เครื่องตีน้ำด้วยอินเวอร์เตอร์ VFD ผ่านโปรโตคอล Modbus RTU เพื่อประหยัดพลังงานไฟฟ้าได้ถึงร้อยละ 40 การวิเคราะห์ปริมาณอาหารเหลือก้นบ่อบนยอยกอาหารกุ้งด้วยคอมพิวเตอร์วิทัศน์ และระบบเตือนภัยฉุกเฉินตลอด 24 ชั่วโมง
\end{enumerate}

คณะผู้เขียนหวังเป็นอย่างยิ่งว่า ตำราเล่มนี้จะเป็นประโยชน์ต่อนักศึกษา นักวิจัย คณาจารย์ เกษตรกรปราดเปรื่อง (Smart Farmers) และผู้สนใจทั่วไป ในการนำองค์ความรู้ไปประยุกต์สร้างสรรค์นวัตกรรมเพื่อยกระดับการเกษตรแม่นยำและการเพาะเลี้ยงสัตว์น้ำของประเทศอย่างยั่งยืนสืบไป

\vspace{1.5cm}
\begin{flushright}
    \textbf{ผู้ช่วยศาสตราจารย์ ดร.ชีวะ ทัศนา}\\
    คณะวิทยาศาสตร์และเทคโนโลยี มหาวิทยาลัยราชภัฏรำไพพรรณี\\
    กันยายน 2570
\end{flushright}
\clearpage
